\section{Native handle feet and variations on the selected physical disk}
\label{sec:ed03-native-feet-jets}

\begin{theorem}[Parametrizations on the actual incoming level]
\label{thm:ed03-native-foot-charts}
\uses{thm:ed03-end-filtrations}
\lean{SmoothPoincare.IsolatedOneHandle.footPartialDiffeomorph, SmoothPoincare.IsolatedOneHandle.footPartialDiffeomorph_apply_nativeAttaching, SmoothPoincare.IsolatedOneHandle.footPartialDiffeomorph_disjoint, SmoothPoincare.IsolatedOneHandle.footComponentChart_closedBall_two_subset_source, SmoothPoincare.IsolatedOneHandle.footComponentPoint_val, SmoothPoincare.IsolatedOneHandle.footComponentChart_symm_point}
\leanok
For each of the two feet of an actual native one-handle in either selected
end filtration, its full unit disk extends to a smooth chart containing the
closed radius-two ball. The two images lie on disjoint branches of the
incoming level. Each chart restricts to the open connected component of
its own foot center. Its inverse recovers the transverse parameter.
\end{theorem}
\begin{proof}
\leanok
In Morse coordinates the transverse coordinate is $rz$ and the signed
axis coordinate has magnitude $\sqrt{2\varepsilon+r^2|z|^2}$. The squared
model radius is $2\varepsilon+2r^2|z|^2$. The retained double-radius
normal form therefore covers the graphs for $|z|\le2$. Division of the
transverse coordinate by $r$ gives the smooth inverse. The image of the
connected closed ball lies in the component of its center.
\end{proof}

\begin{proposition}[Relative oriented feet in a signed collar]
\label{prop:ed03-relative-foot-collar}
\uses{thm:ed03-native-foot-charts}
\lean{SmoothPoincare.IsolatedOneHandle.exists_orientedNativeFootComponentChart, SmoothPoincare.protectedFootBalls_compact_and_compl_preconnected, SmoothPoincare.exists_bufferedFootSignedCollarComparison, SmoothPoincare.SupportedLevelIsotopy.signedCollarComparison_protected, SmoothPoincare.SupportedLevelIsotopy.signedCollarComparison_end}
\leanok
Give the closed native component an orientation. Its complete foot chart
has an oriented representative after either the identity or the central
reflection of the three parameters. A finite pairwise disjoint family of
protected closed chart balls has compact union and preconnected complement.
For a supplied oriented model foot in the same connected level, disjoint
from these protected balls, the supported isotopy matches every foot
parameter. On supplied full signed collar charts it gives a smooth
comparison in both directions. On the central plateau its foot parameter
is the chosen linear isometry; on the protected tubes and both outer
regions it preserves every coordinate under the two markings.
\end{proposition}
\begin{proof}
\leanok
Restrict the actual chart to a connected open source and use the
orientation dichotomy there. Reflection corrects the negative case on
the entire marking. Apply relative ball isotopy off the union of the
protected balls. In signed coordinates the comparison and inverse are
$(x,t)\mapsto(J_{\rho(t)}x,t)$ and $(x,t)\mapsto(J_{\rho(t)}^{-1}x,t)$,
conjugated by the given charts. Choose $\rho=1$ on the central plateau
and $\rho=0$ near both outer ends.
\end{proof}

\begin{theorem}[Calibration by the actual height]
\label{thm:ed03-native-height-collar}
\uses{thm:ed03-end-filtrations}
\lean{SmoothPoincare.exists_nativeSignedHeightCollar, SmoothPoincare.exists_nativeOppositeSignedCollars}
\leanok
On a nonempty regular level $L=h^{-1}(c)$ of a smooth height on a compact
boundaryless source, a smooth partial diffeomorphism with source
$L\times(-r,r)$ has target exactly $h^{-1}((c-r,c+r))$. Its zero slice is
the native inclusion. The inverse signed coordinate at $y$ is $h(y)-c$.
The opposite height has a collar with the same physical image; the
identification is the identity on level points and negates the transverse
coordinate in both directions.
\end{theorem}
\begin{proof}
\leanok
Compactness separates the selected level from the critical values.
Use the smooth gradient-like field with unit negative height speed
throughout a smaller closed regular band. The maps are its flow for time
$-t$ from $x\in L$ and its return flow for time $h(y)-c$ from $y$. The
flow law supplies the two inverse identities. Smooth factorization through
the native level gives smoothness of the inverse projection. Conjugate
by the native opposite-level map and $t\mapsto-t$ to obtain the second chart.
\end{proof}

\begin{proposition}[Physical cutoff jets on the balanced disk]
\label{prop:ed03-physical-marked-jets}
\uses{thm:ed03-atlas-division}
\lean{SmoothPoincare.exists_balancedPhysicalSpectralExteriorDivision_of_endpointData, SmoothPoincare.exists_selectedMarkedRadialCutoff, SmoothPoincare.exists_selectedRadialCutoff_radial_deriv, SmoothPoincare.fderiv_fderiv_selectedRadialCutoff_apply, SmoothPoincare.exists_selectedMarkedMetricLift, SmoothPoincare.selectedMarkedMetricLift_chartGram, SmoothPoincare.exists_selectedRadialSpliceSecondPartialError_on_physicalTrim}
\leanok
Choose one outer marked disk for the two physical chiral measures, the
tested top-eigenvalue measure in the same physical frame, and the full
curvature measure. On a buffered width $w=1/(n+1)$ the physical metric
has a smooth model lift agreeing with its marked pullback on a neighborhood
of the entire closed buffer. The squared-radius step equals one in the
inner region and zero in the outer region. Along every unit ray its
derivative attains $-1/(2w)$ somewhere in the full trim.
For two smooth coefficients $G,B$, write $H=B-G$. In the pullback of the
actual incoming trim volume, the second-partial cutoff error, at path
parameter $0\le t\le1$, has $L^2$ norm bounded by
\[
 t\left[\left(4C_2/w^2+C_1/w\right)\|H\|_2
   +(2C_1/w)\left(\|\partial_m H\|_2+\|\partial_i H\|_2\right)\right].
\]
\end{proposition}
\begin{proof}
\leanok
Transport the restricted volume by the inverse marking. Differentiate
the smooth profile of squared radius twice. The product rule gives one
Hessian term multiplying $H$ and the two displayed mixed terms. The
compact buffer admits a single smooth extension of the pullback metric;
its fixed chart Gram coefficients agree throughout an open neighborhood.
\end{proof}

\begin{theorem}[Attained signed variation on the disk side]
\label{thm:ed03-selected-side-variation}
\uses{prop:ed03-physical-marked-jets}
\lean{SmoothPoincare.coreVariation_taylor_upper, SmoothPoincare.signedCoreSaving_eq_energySaving, SmoothPoincare.positiveCoreVariation_bounds, SmoothPoincare.exists_selectedDiskSideVariation}
\leanok
Put $U=B\cap W_n$ for the chosen outer disk and its marked trim. Test
a two-sided metric path with energy $f(s)=E_+(g_s)$, positive endpoints
at $\pm r$, and the same total charge on the entire complement $U^c$
as the source, separately in each sign. Suppose $f$ is $C^4$ and its
second and fourth derivatives are bounded in absolute value by $D$ on
$[-r,r]$. One of these same two positive endpoints realizes, in each
chirality, a saving on $U$ of at least
\[
 \max\{ |f'(0)|r-Dr^2/2,\;-f''(0)r^2/2-Dr^4/24\}.
\]
The saving is at most the smaller of $f(0)-I$ and the raw charge of $U$
in that sign. For a positive source this yields
$|f'(0)|\le (f(0)-I)/r+Dr/2$ and
$[-f''(0)]_+\le2(f(0)-I)/r^2+Dr^2/12$.
\end{theorem}
\begin{proof}
\leanok
Use Taylor's theorem on both signed intervals and choose the endpoint of
smaller energy. Averaging the cubic Taylor polynomials cancels both odd
terms. Signature and equality on the full complement identify each local
saving with the total drop. Positivity puts that attained endpoint above
the fixed-type infimum. The coordinate width $w$ and variation step $r$
are independent parameters.
\end{proof}

\begin{theorem}[Localized Ricci charge in the modified flow gauge]
\label{thm:ed03-modified-euler-transport}
\lean{SmoothPoincare.metricFrameTensorNormSq_eq_chiral_ricci_scalar, SmoothPoincare.modifiedCutoffMoments_full_eq, SmoothPoincare.modifiedEuler_timeWeight, SmoothPoincare.modifiedEuler_weightedTimeBound, SmoothPoincare.modifiedEuler_transport_of_flux, SmoothPoincare.modifiedTerminalRicci_lower_of_eulerTransport_of_measurableCharges}
\leanok
Let a common cutoff $0\le\chi\le1$ equal one on the initial charge set
and vanish off $B$ at terminal time $\tau>0$. Suppose the initial full
curvature charge is at least $\epsilon^2$. Write $m_+,m_-,m_S$ for
the two cutoff Weyl-square moments and the cutoff scalar-square moment at
time zero. Assume the integrability and measurability of the tested
densities. Given the spatial Euler transgression estimate with coefficient
$K$, a nonnegative remaining time factor of $L^2$ norm at most $L$, and
$b<1/2$, the actual terminal metric satisfies
\[
 \int_B |Z_\tau|^2\,dV_\tau\ \ge \epsilon^2/2-4(m_++m_-)-m_S/6-2KL\sqrt{\tau/(1-2b)}.
\]
\end{theorem}
\begin{proof}
\leanok
The Levi--Civita curvature symmetries and the first Bianchi identity
give the ordered-tensor norm $4(|W^+|^2+|W^-|^2)+2|Z|^2+R^2/6$.
Integrate with the same test on the initial slice. Cauchy--Schwarz in
time uses the exact integral $\tau/(1-2b)$ for the squared time
weight. At the terminal slice drop the nonnegative Weyl and scalar
moments and bound the cutoff Ricci integral by the whole supporting set.
\end{proof}
